\chapter{Obsessive-Compulsive Disorder (OCD)}

\medskip
\noindent\textit{\textbf{Under review.} This chapter is an AI-assisted educational draft; it is not a substitute for professional medical advice.}

\section*{Definition and Overview}
Obsessive-compulsive disorder involves recurrent intrusive thoughts, images, or urges (obsessions) and repetitive behaviours or mental acts (compulsions) performed to reduce distress. The symptoms are time-consuming or cause significant impairment and are not simply preferences or habits.

\section*{Clinical Features}
Common themes include contamination, harm, doubt, order, or taboo thoughts. Compulsions may include washing, checking, counting, repeating, reassurance seeking, or mental rituals. People usually recognise that the cycle is excessive, although insight varies.

\section*{Diagnosis and Differential Diagnosis}
Diagnosis is clinical and considers duration, distress, impairment, insight, tics, depression, anxiety, psychosis, autism, substance effects, and medical causes. Screening tools may support but cannot replace assessment.

\section*{Management}
First-line treatment is cognitive behavioural therapy with exposure and response prevention. Selective serotonin reuptake inhibitors may be prescribed when indicated; severe or persistent symptoms may need specialist care. Reassurance and avoidance can unintentionally maintain the cycle.

\section*{When to Seek Help}
Seek assessment when symptoms interfere with daily life, relationships, study, or work. Seek urgent help for self-harm thoughts, inability to eat or drink, or inability to remain safe.

\section*{Self-Care}
Follow the treatment plan, practise agreed exposure exercises, maintain sleep and routines, and involve trusted supports with consent. Do not stop prescribed medicines abruptly.

\section*{Expanded clinical context}
OCD is a mental-health topic. Interpretation should account for age, baseline health, medicines, comorbidities, goals, and access to care. This chapter supports discussion with a qualified clinician and does not establish a diagnosis.

\section*{Assessment and decision points}
Assessment should consider onset, duration, triggers, sleep, substance use, physical health, developmental context, social stressors, and immediate safety. Ask about self-harm or harm from others when appropriate. Document the main concern, time course, functional impact, relevant risks, and red flags. Use targeted investigations when their result is likely to change management.

\section*{Management and follow-up}
Management may include education, psychological therapy, practical support, treatment of coexisting conditions, medication when indicated, and an agreed follow-up or safety plan. Explain expected improvement, when reassessment is needed, and how follow-up can be accessed. Avoid starting, stopping, or sharing prescription medicines without professional advice.

\section*{Safety-netting}
Seek urgent help for immediate danger, suicidal intent, psychosis, severe confusion, inability to meet basic needs, or rapidly worsening symptoms.

\section*{Practical prevention}
Use plain-language information, supportive relationships, regular routines, and professional review when symptoms persist or impair daily life. Keep written instructions and seek review if the topic recurs, changes, or does not improve as expected.

\section*{Limitations}
This educational expansion should be checked against current local guidance and authoritative topic-specific sources.
\clearpage
